% ============================================================
%  PV DE RECETTE FONCTIONNELLE — Phase 5 : Clôture
%  Time'Eats — Cellule PMO
% ============================================================
\documentclass[11pt,a4paper,oneside]{report}
\usepackage{../style/consulting}
\hypersetup{pdftitle={PV de Recette Fonctionnelle — Time'Eats PMO}}
\pagestyle{fancy}\fancyhf{}
\renewcommand{\headrulewidth}{0pt}
\fancyhead[L]{\begin{tikzpicture}[remember picture,overlay]
  \draw[cnNavy,line width=0.5pt](0,0)--(\textwidth,0);
\end{tikzpicture}\vspace{2pt}\timeeatslogo\hspace{4pt}\small\textcolor{cnSlate}{Time'Eats \textbf{·} PV de Recette Fonctionnelle — \textit{[Nom projet]}}}
\fancyhead[R]{\small\textcolor{cnSlate}{\thepage}}
\fancyfoot[C]{\tiny\textcolor{cnSlate}{Document confidentiel — Cellule PMO — CESI MAALSI 2026}}
\begin{document}\small

\begin{center}
{\LARGE\bfseries\color{cnNavy} Procès-Verbal de Recette Fonctionnelle}\\[4pt]
\textcolor{cnOrange}{\rule{10cm}{1pt}}\\[4pt]
{\large\color{cnSlate} Phase 5 — Clôture \quad|\quad Projet : \textit{[Nom du projet]}}
\end{center}

\vspace{4pt}
\renewcommand{\arraystretch}{1.4}
\begin{tabularx}{\textwidth}{L{3.5cm} Y L{3.5cm} Y}
\toprule
\rowcolor{cnNavy}\th{Champ} & \th{Valeur} & \th{Champ} & \th{Valeur} \\
\midrule
Projet          & \textit{[Nom]}           & Version livrée  & \textit{[v X.X]} \\
\rowcolor{cnIce}
Date de recette & \textit{[JJ/MM/AAAA]}    & Env. de recette & \textit{[Préprod / UAT]} \\
Chef de projet  & \textit{[Prénom NOM]}    & Responsable MOA & \textit{[Prénom NOM]} \\
\bottomrule
\end{tabularx}

\section*{1 — Périmètre testé}

\renewcommand{\arraystretch}{1.3}
\begin{tabularx}{\textwidth}{C{1cm} L{5cm} Y C{3cm}}
\toprule
\rowcolor{cnNavy}\th{\#} & \th{Fonctionnalité / Module} & \th{Critères d'acceptance} & \th{Statut} \\
\midrule
F01 & \textit{[Fonctionnalité 1]} & \textit{[Critère mesurable]} & $\square$ OK \quad $\square$ Réserve \quad $\square$ KO \\
\rowcolor{cnIce}
F02 & \textit{[Fonctionnalité 2]} & \textit{[Critère mesurable]} & $\square$ OK \quad $\square$ Réserve \quad $\square$ KO \\
F03 & \textit{[Fonctionnalité 3]} & \textit{[Critère mesurable]} & $\square$ OK \quad $\square$ Réserve \quad $\square$ KO \\
\rowcolor{cnIce}
F04 & \textit{[Fonctionnalité 4]} & \textit{[Critère mesurable]} & $\square$ OK \quad $\square$ Réserve \quad $\square$ KO \\
F05 & \textit{[Fonctionnalité 5]} & \textit{[Critère mesurable]} & $\square$ OK \quad $\square$ Réserve \quad $\square$ KO \\
\bottomrule
\end{tabularx}

\section*{2 — Résultats globaux}

\renewcommand{\arraystretch}{1.3}
\begin{tabularx}{\textwidth}{L{4cm} C{2cm} C{2cm} C{2cm} Y}
\toprule
\rowcolor{cnNavy}\th{Catégorie} & \th{Total} & \th{OK} & \th{Réserve} & \th{KO} \\
\midrule
Fonctionnalités testées & \textit{[N]} & \textit{[N]} & \textit{[N]} & \textit{[N]} \\
\rowcolor{cnIce}
Anomalies P1 bloquantes & -- & \textit{[N résolu]} & -- & \textit{[N ouvert]} \\
Anomalies P2 majeures   & -- & \textit{[N résolu]} & -- & \textit{[N ouvert]} \\
\bottomrule
\end{tabularx}

\section*{3 — Réserves émises}

\begin{enumerate}
  \item \textit{[Réserve 1 : description + engagement de correction avant \textit{[date]}]}
  \item \textit{[Réserve 2 : description + engagement]}
  \item \textit{[Aucune réserve — le cas échéant]}
\end{enumerate}

\section*{4 — Décision de recette}

\begin{tcolorbox}[
  colback=cnLightBlue, colframe=cnNavy,
  leftrule=4pt, rightrule=0pt, toprule=0pt, bottomrule=0pt,
  arc=0pt, boxsep=6pt, top=6pt, bottom=6pt, left=6pt, right=6pt]

\renewcommand{\arraystretch}{1.5}
\begin{tabularx}{\textwidth}{ccc}
\cellcolor{cnGreen!25}\large $\square$ \textbf{Recette acceptée}
& \cellcolor{cnOrange!25}\large $\square$ \textbf{Acceptée avec réserves}
& \cellcolor{cnRed!25}\large $\square$ \textbf{Recette refusée} \\
\end{tabularx}

\vspace{6pt}
\textbf{Conditions de levée des réserves :} \textit{[Date butoir et engagements ESN]}

\end{tcolorbox}

\section*{5 — Signatures}

\vspace{4pt}
\renewcommand{\arraystretch}{2.5}
\begin{tabularx}{\textwidth}{Y Y Y}
\toprule
\rowcolor{cnNavy}\th{Responsable MOA} & \th{Chef de projet} & \th{DSI} \\
\midrule
\textit{Nom :} & \textit{Nom :} & \textit{Nom :} \\
\textit{Fonction :} & \textit{Fonction :} & \textit{Fonction :} \\
\textit{Date :} & \textit{Date :} & \textit{Date :} \\
\textit{Signature :} & \textit{Signature :} & \textit{Signature :} \\
\bottomrule
\end{tabularx}

\end{document}
